%============================================================
% Research Progress Report — Weeks 7 & 8
% Reduction Ladder for Code & Multi-Arm Mitigation
% Orange Innovation Labs — AI Research & Development Division
% Authors: Omar Abdelhamid & Nour Walid (Equal Contribution)
% Submitted to: Dr. Ghada Soliman
% Compile with: pdflatex -> bibtex -> pdflatex x2
%============================================================
\documentclass[12pt, a4paper]{article}

%------------------------------------------------------------
% PACKAGES
%------------------------------------------------------------
\usepackage[T1]{fontenc}
\usepackage[utf8]{inputenc}
\usepackage{lmodern}
\usepackage{microtype}
\usepackage[left=2.2cm, right=2.2cm, top=2.8cm, bottom=2.8cm]{geometry}
\usepackage{xcolor}
\usepackage{graphicx}
\usepackage{amsmath, amssymb}
\usepackage{booktabs}
\usepackage{tabularx}
\usepackage{longtable}
\usepackage{array}
\usepackage{multirow}
\usepackage{fancyhdr}
\usepackage{titlesec}
\usepackage{tocloft}
\usepackage{hyperref}
\usepackage{enumitem}
\usepackage{mdframed}
\usepackage{tikz}
\usepackage{pgfplots}
\pgfplotsset{compat=1.18}
\usepackage{float}
\usepackage{caption}
\usepackage{subcaption}
\usepackage{setspace}
\allowdisplaybreaks
\usepackage{parskip}
\usepackage{soul}
\usepackage{pifont}
\usepackage{fontawesome5}
\usepackage{listings}
\usepackage{tcolorbox}
\tcbuselibrary{skins, breakable, listings}
\usepackage[backend=bibtex, bibstyle=ieee, citestyle=numeric-comp, sorting=none, maxbibnames=4]{biblatex}
\addbibresource{references.bib}

%------------------------------------------------------------
% ORANGE BRAND COLORS
%------------------------------------------------------------
\definecolor{OrangeMain}{RGB}{255, 100, 0}
\definecolor{OrangeDark}{RGB}{200, 70, 0}
\definecolor{OrangeLight}{RGB}{255, 165, 60}
\definecolor{DarkBG}{RGB}{20, 20, 25}
\definecolor{MidGray}{RGB}{80, 80, 90}
\definecolor{LightGray}{RGB}{240, 240, 245}
\definecolor{AccentBlue}{RGB}{30, 120, 200}
\definecolor{SuccessGreen}{RGB}{40, 160, 80}
\definecolor{AlertRed}{RGB}{200, 40, 40}
\definecolor{Week8Color}{RGB}{100, 50, 200}

%------------------------------------------------------------
% HYPERREF SETUP
%------------------------------------------------------------
\hypersetup{
  colorlinks   = true,
  linkcolor    = OrangeDark,
  citecolor    = OrangeDark,
  urlcolor     = AccentBlue,
  pdftitle     = {Weeks 7--8 Progress Report: Reduction Ladder for Code \& Multi-Arm Mitigation},
  pdfauthor    = {Omar Abdelhamid, Nour Walid},
  pdfsubject   = {Bi-Weekly Research Progress Report for Dr. Ghada Soliman - Orange Innovation Labs},
}

%------------------------------------------------------------
% SECTION HEADING STYLE
%------------------------------------------------------------
\titleformat{\section}
  {\large\bfseries\color{OrangeMain}}
  {\thesection.}{0.6em}{}
  [\color{OrangeMain}\titlerule]

\titleformat{\subsection}
  {\normalsize\bfseries\color{OrangeDark}}
  {\thesubsection.}{0.5em}{}

\titleformat{\subsubsection}
  {\normalsize\itshape\color{MidGray}}
  {\thesubsubsection.}{0.4em}{}

%------------------------------------------------------------
% HEADER / FOOTER
%------------------------------------------------------------
\pagestyle{fancy}
\fancyhf{}
\renewcommand{\headrulewidth}{0.5pt}
\renewcommand{\headrule}{\hbox to\headwidth{\color{OrangeMain}\leaders\hrule height\headrulewidth\hfill}}
\fancyhead[L]{\small\color{MidGray} Weeks 7--8 Research Report \textbar\ Reduction Ladder \& Mitigation}
\fancyhead[R]{\small\color{MidGray} Orange Innovation Labs}
\fancyfoot[C]{\small\color{MidGray} \thepage}
\fancyfoot[R]{\small\color{OrangeMain} \textbf{Confidential --- Orange Restricted}}

%------------------------------------------------------------
% CUSTOM BOXES
%------------------------------------------------------------
\newtcolorbox{researchbox}[1][]{
  enhanced, breakable,
  colback=LightGray,
  colframe=OrangeMain,
  fonttitle=\bfseries\color{white},
  coltitle=white,
  attach boxed title to top left={yshift=-2mm, xshift=4mm},
  boxed title style={colback=OrangeMain, rounded corners},
  arc=3pt, boxrule=1pt,
  #1
}

\newtcolorbox{domainbox}[1][]{
  enhanced, breakable,
  colback=AccentBlue!6,
  colframe=AccentBlue,
  fonttitle=\bfseries\color{white},
  coltitle=white,
  attach boxed title to top left={yshift=-2mm, xshift=4mm},
  boxed title style={colback=AccentBlue, rounded corners},
  arc=3pt, boxrule=1pt,
  #1
}

\newtcolorbox{week8box}[1][]{
  enhanced, breakable,
  colback=Week8Color!5,
  colframe=Week8Color,
  fonttitle=\bfseries\color{white},
  coltitle=white,
  attach boxed title to top left={yshift=-2mm, xshift=4mm},
  boxed title style={colback=Week8Color, rounded corners},
  arc=3pt, boxrule=1pt,
  #1
}

\newtcolorbox{gapbox}{
  enhanced,
  colback=AlertRed!8,
  colframe=AlertRed,
  arc=3pt, boxrule=0.8pt,
  fonttitle=\bfseries,
  title={\faExclamationTriangle\ Scientific Integrity Alert: Test-Set Contamination},
  coltitle=AlertRed,
  attach boxed title to top left={yshift=-2mm, xshift=4mm},
  boxed title style={colback=white, colframe=AlertRed},
}

\newtcolorbox{noveltybox}{
  enhanced,
  colback=SuccessGreen!8,
  colframe=SuccessGreen,
  arc=3pt, boxrule=0.8pt,
  fonttitle=\bfseries,
  title={\faCheckCircle\ Primary Innovation: Invariance-Regularized Policy Optimization (Inv-GRPO)},
  coltitle=SuccessGreen,
  attach boxed title to top left={yshift=-2mm, xshift=4mm},
  boxed title style={colback=white, colframe=SuccessGreen},
}

\newtcolorbox{keybox}[1][]{
  enhanced, breakable,
  colback=DarkBG!5,
  colframe=OrangeDark,
  left=4pt, right=4pt,
  arc=2pt, boxrule=0.8pt,
  #1
}

\newtcolorbox{findingbox}{
  enhanced,
  colback=OrangeLight!10,
  colframe=OrangeMain,
  arc=3pt, boxrule=0.8pt,
  fonttitle=\bfseries,
  title={\faLightbulb\ Key Empirical Finding: Vanilla SFT Trade-off Revealed},
  coltitle=OrangeDark,
  attach boxed title to top left={yshift=-2mm, xshift=4mm},
  boxed title style={colback=white, colframe=OrangeMain},
}

%------------------------------------------------------------
% LISTINGS (code style)
%------------------------------------------------------------
\lstdefinestyle{pythonstyle}{
  language=Python,
  basicstyle=\ttfamily\footnotesize,
  keywordstyle=\color{OrangeMain}\bfseries,
  commentstyle=\color{MidGray}\itshape,
  stringstyle=\color{SuccessGreen},
  numberstyle=\tiny\color{MidGray},
  numbers=left,
  stepnumber=1,
  breaklines=true,
  frame=single,
  rulecolor=\color{LightGray},
  backgroundcolor=\color{LightGray},
}

%============================================================
% DOCUMENT BEGIN
%============================================================
\begin{document}

%------------------------------------------------------------
% COVER PAGE
%------------------------------------------------------------
\begin{titlepage}
  \pagecolor{DarkBG}
  \color{white}
  \begin{center}
    \vspace*{0.8cm}
    
    {\color{OrangeMain}\rule{\textwidth}{3pt}}
    \vspace{0.4cm}
    
    {\Large\bfseries\color{OrangeLight} Orange Innovation Labs}\\[0.2cm]
    {\normalsize\color{MidGray} AI Research \& Development Division --- Egypt}
    
    \vspace{0.6cm}
    {\color{OrangeMain}\rule{0.4\textwidth}{1pt}}
    \vspace{0.6cm}
    
    {\Huge\bfseries\color{white}
      Bi-Weekly Research Progress Report\\[0.3cm]
      \color{OrangeMain} Weeks 7 \& 8: From Baseline Diagnostics\\
      \color{white} to Arm~1 Inv-GRPO Training \& First Results}\\[0.4cm]
    
    {\large\color{LightGray}
      Comprehensive Account of Stage 3 Distillation, QLoRA Fine-Tuning of M2,\\
      Post-Training Comparative Evaluation (M1 vs.\ M2), the Four Mitigation Arms,\\
      and the Full Implementation of Invariance-Regularized Policy Optimization (Inv-GRPO)}
    
    \vspace{0.8cm}
    {\color{OrangeMain}\rule{0.4\textwidth}{1pt}}
    \vspace{0.8cm}
    
    \begin{tabular}{rl}
      \textbf{\color{OrangeLight} Research Authors:}   & \textbf{Omar Abdelhamid} \& \textbf{Nour Walid}\\[2pt]
                                                       & {\small AI R\&D, Orange Innovation Egypt \textbar\ Collaborative Authorship}\\[6pt]
      \textbf{\color{OrangeLight} Submitted to:}       & \textbf{Dr.\ Ghada Soliman}\\[2pt]
                                                       & {\small Head of Software Engineering \& AI Research, Orange Innovation Labs}\\[6pt]
      \textbf{\color{OrangeLight} Reporting Period:}   & \textbf{Week 7} (Post-Baseline: Distillation \& Infrastructure) +\\
                                                       & \textbf{Week 8} (QLoRA Training, M2 Eval, Arm 1 Inv-GRPO)\\[4pt]
      \textbf{\color{OrangeLight} Primary Target:}     & \texttt{Qwen/Qwen2.5-Coder-1.5B-Instruct} (Edge SLM)\\[4pt]
      \textbf{\color{OrangeLight} Hardware Node:}      & NVIDIA RTX 3070 Ti Laptop (8 GB GDDR6 VRAM, CC~8.6)\\[4pt]
      \textbf{\color{OrangeLight} Models Evaluated:}   & \textbf{M1} (Zero-Shot Baseline), \textbf{M2} (Vanilla SFT QLoRA)\\[4pt]
      \textbf{\color{OrangeLight} Report Status:}      & \textcolor{OrangeLight}{\bfseries Final Submission --- Week 8 Milestone}\\[4pt]
      \textbf{\color{OrangeLight} Security Level:}     & \textcolor{OrangeMain}{\bfseries Orange Restricted / Confidential}\\
    \end{tabular}
    
    \vfill
    {\color{OrangeMain}\rule{\textwidth}{3pt}}
    \vspace{0.2cm}
    {\small\color{MidGray}
      Proprietary Technical Document prepared for Orange Innovation Labs AI R\&D Division.}
  \end{center}
\end{titlepage}

\nopagecolor

%------------------------------------------------------------
% TABLE OF CONTENTS
%------------------------------------------------------------
\tableofcontents
\newpage

%============================================================
% SECTION 1: EXECUTIVE SUMMARY
%============================================================
\section{Executive Summary: Two-Week Consolidated Overview}

\begin{researchbox}[title={\faStar\ Two-Week Research Snapshot}]
\begin{itemize}[leftmargin=*]
  \item \textbf{5 Notebooks Completed \& Verified:} \texttt{nb\_01} through \texttt{nb\_05}, covering the full pipeline from data acquisition to post-training comparative evaluation.
  \item \textbf{2 Models Empirically Evaluated:} M1 (zero-shot baseline) and M2 (Vanilla SFT via QLoRA), measured across \textbf{664 benchmark tasks} spanning 6 Reduction Ladder levels ($L_0 \dots L_5$) \cite{evoeval}.
  \item \textbf{Key Finding:} Vanilla SFT (M2) improves Ladder AUC by $+1.05$\,pp ($82.22\% \to 83.27\%$) and dramatically boosts robustness on higher rungs ($L_2 \dots L_5$, up to $+18$\,pp), but exhibits a canonical SFT trade-off: a $-28.66\%$ drop on $L_0$ canonical HumanEval tasks.
  \item \textbf{Arm 1 (Inv-GRPO) Fully Implemented:} The \texttt{InvGRPOTrainer} is memory-optimized with micro-batching ($\approx 3.03$\,GB peak VRAM), resolving OOM on RTX 3070 Ti.
  \item \textbf{4 Research Arms Engineered:} All four mitigation notebooks complete with verified reward engines.
  \item \textbf{9 Bugs Resolved:} Full post-mortem covering TRL 0.24, Transformers 4.51, Pydantic, Sandbox, and VRAM.
\end{itemize}
\end{researchbox}

\vspace{0.4cm}
The current milestone establishes a rigorous, contamination-free empirical foundation, revealing both the gains and costs of naive chain-of-thought distillation. This motivates our principal contribution---\textbf{Inv-GRPO}---which targets higher-rung robustness \emph{without} sacrificing $L_0$ canonical coding precision.

%============================================================
% SECTION 2: WEEK 7 WORK
%============================================================
\section{Week 7: Infrastructure, Distillation, \& Mitigation Arms Engineering}

\subsection{Notebook NB-01: Data Pipeline \& Benchmark Acquisition}
\textbf{Notebook:} \texttt{notebooks/nb\_01\_data\_pipeline.ipynb}

\begin{researchbox}[title={\faDatabase\ NB-01 Acquisition Statistics}]
\begin{itemize}[leftmargin=*]
  \item \textbf{HumanEval ($L_0$):} 164 canonical coding problems cached locally.
  \item \textbf{EvoEval ($L_1$--$L_5$) \cite{evoeval}:} 100 tasks/level $\times$ 5 levels = \textbf{500 held-out evaluation tasks}.
  \item \textbf{Total Benchmark Suite:} \textbf{664 tasks} across 6 semantic perturbation levels.
  \item \textbf{Cache Strategy:} \texttt{src/infrastructure/hf\_loader.py} enforces offline-first Arrow-backed caching---no HuggingFace Hub dependency during training loops.
  \item \textbf{Ground-Truth Validation:} All assertion suites independently verified via SubprocessSandbox \cite{livecodebench}.
\end{itemize}
\end{researchbox}

\subsection{Notebook NB-02: Baseline Evaluation (Model M1)}
\textbf{Notebook:} \texttt{notebooks/nb\_02\_baseline\_eval.ipynb}\\
\textbf{Reference Literature:} EvoEval \cite{evoeval}, GSM-Symbolic \cite{gsm_symbolic}, Memorize-or-Generalize \cite{memorize_or_generalize}.

Full zero-shot benchmark sweep of \texttt{Qwen2.5-Coder-1.5B-Instruct} (M1) across 664 tasks.
Results persisted at \texttt{results/baseline/suite\_report\_baseline.json} ($21.3$\,MB).

\begin{table}[H]
\centering
\small
\caption{Model M1 (Zero-Shot Baseline)---Per-Level Pass@1 on the 6-Level Reduction Ladder (from \texttt{nb\_02\_baseline\_eval.ipynb}).}
\label{tab:m1_per_level}
\begin{tabularx}{\textwidth}{c l r r r}
\toprule
\textbf{Rung} & \textbf{Perturbation Type} & \textbf{Tasks} & \textbf{Pass@1} & \textbf{$\Delta$ from $L_0$} \\
\midrule
$L_0$ & Canonical HumanEval (no perturbation) \cite{evoeval}  & 164 & \textbf{93.29\%} & --- \\
$L_1$ & Subtle renaming \& docstring reframing               & 100 & 87.00\%          & $-6.29$\,pp \\
$L_2$ & Verbose frame shift / API expansion                  & 100 & 83.00\%          & $-10.29$\,pp \\
$L_3$ & Semantic recomposition                               & 100 & 77.00\%          & $-16.29$\,pp \\
$L_4$ & Algorithmic analogue shift                           & 100 & 77.00\%          & $-16.29$\,pp \\
$L_5$ & Deep structural inversion                            & 100 & 76.00\%          & $-17.29$\,pp \\
\midrule
\multicolumn{2}{l}{\textbf{Ladder AUC} $\mathcal{A}_{M1}$}      & 664 & \multicolumn{2}{c}{\textbf{82.22\%}} \\
\multicolumn{2}{l}{\textbf{Collapse Point} $\ell^*$}            & ---  & \multicolumn{2}{c}{None (all rungs $>50\%$)} \\
\multicolumn{2}{l}{\textbf{Consistency Delta} $\Delta C$}       & ---  & \multicolumn{2}{c}{$0.040$ (4.0\%)} \\
\multicolumn{2}{l}{\textbf{Memorization Risk Index} $\mathcal{M}$} & --- & \multicolumn{2}{c}{$0.163$ (moderate)} \\
\bottomrule
\end{tabularx}
\end{table}

\noindent The M1 baseline confirms our foundational hypothesis: a clear monotone degradation trend \cite{evoeval, gsm_symbolic} from 93.29\% at $L_0$ to 76.00\% at $L_5$ ($-17.29$\,pp). The Memorization Risk of 0.163 indicates a non-trivial fraction of $L_0$ performance derives from surface pattern memorization rather than generalizable reasoning \cite{memorize_or_generalize, learned_or_memorized}.

\subsection{Notebook NB-03: Stage 3 Distillation---Vanilla CoT Corpus}
\textbf{Notebook:} \texttt{notebooks/nb\_03\_distillation.ipynb}\\
\textbf{Reference Literature:} SuperCorrect \cite{supercorrect}, DeepSeekMath GRPO \cite{grpo_deepseek}.

\begin{researchbox}[title={\faDatabase\ NB-03 Distillation Statistics}]
\begin{itemize}[leftmargin=*]
  \item \textbf{Output:} \texttt{data/distillation/sft\_positive\_cot.jsonl}
  \item \textbf{Total Verified Samples:} $\mathbf{500}$ algorithmic CoT traces
  \item \textbf{Token Distribution:} Min $162$ / Max $895$ / \textbf{Mean $377$ tokens/sample}
  \item \textbf{Context Safety:} $100\%$ of samples fit within the 2048-token limit of Qwen2.5-Coder \cite{qwen25coder}.
  \item \textbf{Thought Format:} 3-stage hierarchical template (Problem Decomposition $\to$ Strategy $\to$ Code), conforming to SuperCorrect \cite{supercorrect}.
  \item \textbf{Distribution Plot:} \texttt{results/vanilla\_token\_dist.png}.
\end{itemize}
\end{researchbox}

\subsubsection{Contamination Discovery \& Decoupling}

\begin{gapbox}
\texttt{contrastive\_builder.py} extracted training pairs directly from EvoEval ($L_1$--$L_5$), constituting \textbf{direct benchmark leakage}: the model would memorize test perturbations rather than learn invariant reasoning.
\end{gapbox}

\begin{noveltybox}
\textbf{Resolution:} (1) Purified Stage 3---all EvoEval references removed. (2) Decoupled contrastive DPO to \texttt{arm\_02}, sourced from SuperCorrect \cite{supercorrect} and ReCode \cite{recode_contrastive}. (3) Full Reduction Ladder (664 tasks) remains \textbf{100\% held-out} from all training.
\end{noveltybox}

\subsection{Multi-Arm Mitigation Framework (Week 7)}

\begin{table}[H]
\centering
\small
\caption{Taxonomy of Implemented Multi-Arm Mitigation Approaches.}
\label{tab:mitigation_arms}
\begin{tabularx}{\textwidth}{c l p{3.6cm} p{3.4cm} c}
\toprule
\textbf{Arm} & \textbf{Paradigm} & \textbf{Mechanism} & \textbf{Literature} & \textbf{Overhead} \\
\midrule
\textbf{1} & \textbf{Inv-GRPO}       & Paired invariance reward + template penalty & Primary contribution         & \textbf{0 ms} \\
\textbf{2} & \textbf{Contrastive DPO} & $y^+$ vs $y^-$ shortcut rejection          & \cite{supercorrect, recode_contrastive} & 0 ms \\
\textbf{3} & \textbf{AST-RL}          & Syntax-tree Jaccard distance reward         & \cite{treediff, veriseek}              & 0 ms \\
\textbf{4} & \textbf{Step-RLVR}       & Stepwise contract-gated execution           & \cite{codeprm, execverify}             & +PRM \\
\bottomrule
\end{tabularx}
\end{table}

\textbf{Arm 3 Verified:} Variable renaming preserves $>95\%$ AST similarity; structurally wrong code drops to $<20\%$.\\
\textbf{Arm 4 Verified:} 4/5-step partially-correct solution retains \textbf{0.70 partial credits} vs.\ \textbf{0.00} under binary RLVR.

%============================================================
% SECTION 3: WEEK 8 WORK
%============================================================
\section{Week 8: QLoRA Training, Post-Training Evaluation, \& Arm 1 Inv-GRPO}

\begin{week8box}[title={\faFlask\ Week 8 Milestones}]
\begin{enumerate}[leftmargin=*, label=\textbf{\arabic*.}]
  \item Executed QLoRA fine-tuning of M2 via \texttt{nb\_04\_qlora\_training.ipynb}.
  \item Completed M1 vs.\ M2 comparative evaluation via \texttt{nb\_05\_post\_training\_eval.ipynb} (664 tasks).
  \item Implemented full \texttt{InvGRPOTrainer} with micro-batched backward pass ($\approx 3.03$\,GB VRAM).
  \item Resolved 3 additional infrastructure bugs (BUG-7, BUG-8, BUG-9).
  \item Pushed complete modular codebase to GitHub.
\end{enumerate}
\end{week8box}

\subsection{Notebook NB-04: QLoRA Fine-Tuning of Model M2}
\textbf{Notebook:} \texttt{notebooks/nb\_04\_qlora\_training.ipynb}\\
\textbf{Reference Literature:} QLoRA \cite{peft_qlora}, TRL SFTTrainer.

\begin{table}[H]
\centering
\small
\caption{QLoRA Fine-Tuning Configuration for Model M2 (Vanilla SFT).}
\label{tab:qlora_config}
\begin{tabularx}{\textwidth}{l l X}
\toprule
\textbf{Parameter} & \textbf{Value} & \textbf{Rationale} \\
\midrule
Base Model           & \texttt{Qwen2.5-Coder-1.5B-Instruct} & Target edge SLM \cite{qwen25coder} \\
Quantization         & 4-bit NF4 (BitsAndBytes 0.49.2)      & 1.12 GB base VRAM footprint \\
LoRA Rank $r$        & $16$                                  & Efficiency-capacity balance \cite{peft_qlora} \\
LoRA $\alpha$        & $32$                                  & Standard $\alpha = 2r$ scaling \\
Target Modules       & \texttt{q, k, v, o\_proj}             & All attention projections \\
Max Sequence Length  & $2048$                                & Covers $100\%$ corpus samples \\
Training Samples     & $500$                                 & Verified CoT traces \\
Peak VRAM (Training) & $6.42$\,GB / $8.00$\,GB               & $80.2\%$ utilization (safe budget) \\
Grad Checkpointing   & Enabled                               & Halves activation memory \\
\bottomrule
\end{tabularx}
\end{table}

\subsection{Notebook NB-05: Post-Training Comparative Evaluation (M1 vs.\ M2)}
\textbf{Notebook:} \texttt{notebooks/nb\_05\_post\_training\_eval.ipynb}\\
\textbf{Reference Literature:} EvoEval \cite{evoeval}, LiveCodeBench \cite{livecodebench}, Flip-Flop Consistency \cite{flip_flop_consistency}, IB-FT \cite{ibft_memorization}.\\
\textbf{Persisted Results:} \texttt{results/distilled\_vanilla/suite\_report\_distilled\_vanilla.json} ($20.9$\,MB).

\begin{table}[H]
\centering
\small
\caption{Empirical Per-Level Pass@1: M1 vs.\ M2 (from \texttt{nb\_05\_post\_training\_eval.ipynb}, 664 held-out tasks).}
\label{tab:m1_vs_m2}
\begin{tabularx}{\textwidth}{c l r r r r}
\toprule
\textbf{Rung} & \textbf{Perturbation Type} & \textbf{Tasks} & \textbf{M1 Pass@1} & \textbf{M2 Pass@1} & \textbf{$\Delta$ M2$-$M1} \\
\midrule
$L_0$ & Canonical HumanEval \cite{evoeval}     & 164 & 93.29\% & 64.63\% & \textcolor{AlertRed}{$-28.66$\,pp} \\
$L_1$ & Subtle renaming / docstring reframing   & 100 & 87.00\% & 60.00\% & \textcolor{AlertRed}{$-27.00$\,pp} \\
$L_2$ & Verbose frame shift / API expansion     & 100 & 83.00\% & \textbf{96.00\%} & \textcolor{SuccessGreen}{$+13.00$\,pp} \\
$L_3$ & Semantic recomposition                  & 100 & 77.00\% & \textbf{95.00\%} & \textcolor{SuccessGreen}{$+18.00$\,pp} \\
$L_4$ & Algorithmic analogue shift              & 100 & 77.00\% & \textbf{93.00\%} & \textcolor{SuccessGreen}{$+16.00$\,pp} \\
$L_5$ & Deep structural inversion               & 100 & 76.00\% & \textbf{91.00\%} & \textcolor{SuccessGreen}{$+15.00$\,pp} \\
\midrule
\multicolumn{3}{l}{\textbf{Ladder AUC} $\mathcal{A}$}         & 82.22\% & 83.27\% & \textcolor{SuccessGreen}{$+1.05$\,pp} \\
\multicolumn{3}{l}{\textbf{Consistency Delta} $\Delta C$}      & 0.040   & 0.360   & \textcolor{AlertRed}{$\uparrow$ (less consistent)} \\
\multicolumn{3}{l}{\textbf{Memorization Risk} $\mathcal{M}$}  & 0.163   & \textbf{0.000}   & \textcolor{SuccessGreen}{$-0.163$ (fully resolved)} \\
\bottomrule
\end{tabularx}
\end{table}

\begin{findingbox}
\textbf{Central Empirical Finding (NB-05):} Vanilla SFT (M2) produces a bifurcated performance profile: dramatic gains on structurally perturbed rungs ($L_2{:}+13$\,pp, $L_3{:}+18$\,pp, $L_4{:}+16$\,pp, $L_5{:}+15$\,pp) paired with substantial degradation on canonical tasks ($L_0{:}-28.7$\,pp, $L_1{:}-27.0$\,pp). Memorization Risk drops to exactly \textbf{0.000}, confirming M2 no longer pattern-matches---but at the cost of general coding fluency \cite{ibft_memorization, supercorrect}. This bifurcation constitutes the empirical motivation for Inv-GRPO: achieving M2-level structural robustness \emph{without} sacrificing M1-level $L_0$ precision.
\end{findingbox}

\subsection{Arm 1: Invariance-Regularized Policy Optimization (Inv-GRPO)}
\textbf{Notebook:} \texttt{notebooks/arm\_01\_inv\_grpo.ipynb}\\
\textbf{Implementation:} \texttt{src/arms/arm1\_inv\_grpo/trainer.py}\\
\textbf{Reference Literature:} DeepSeekMath GRPO \cite{grpo_deepseek}, CLARity \cite{clarity_rl}, LoPE \cite{lope_perturbation}, Break-The-Chain \cite{break_the_chain}, CLIPO \cite{clipo_rl}.

\subsubsection{Modified GRPO Objective}

\begin{equation}
\mathcal{R}_{\text{inv}}(y_i, y'_i) = \underbrace{\mathcal{R}_{\text{exec}}(y_i) + \mathcal{R}_{\text{exec}}(y'_i)}_{\text{dual execution}} + \underbrace{\lambda \cdot \mathbb{I}(\text{both pass})}_{\text{consistency bonus}} - \underbrace{\gamma \cdot \mathcal{P}_{\text{template}}(y'_i, y_{\text{decoy}})}_{\text{shortcut penalty}}
\end{equation}

\subsubsection{VRAM-Optimized Micro-Batched Implementation}

\begin{week8box}[title={\faMemory\ OOM Resolution: Micro-Batched Inv-GRPO}]
The naive GRPO forward pass materializes a $(G \times T \times V)$ logit tensor: for $G=8$, $T=512$, $V=151{,}936$ vocabulary, this requires $\approx 12.3$\,GB VRAM---exceeding our 8\,GB budget by $54\%$ \cite{grpo_deepseek}.

\textbf{Solution:} \texttt{InvGRPOTrainer.\_backward\_group\_loss()} processes each $(x_i, x'_i)$ pair sequentially in micro-batches, accumulating gradients without materializing the full tensor. Results:
\begin{itemize}[leftmargin=*]
  \item \textbf{Na\"ive GRPO:} $\approx 12.3$\,GB (OOM crash on RTX 3070 Ti)
  \item \textbf{Micro-Batched Inv-GRPO:} $\approx \mathbf{3.03}$\,GB peak (\textbf{75\% VRAM reduction})
  \item \textbf{Remaining headroom:} $\approx 4.97$\,GB for gradient accumulation
\end{itemize}
\end{week8box}

\subsubsection{Reward Engine Verification}

\begin{table}[H]
\centering
\small
\caption{Inv-GRPO Reward Engine Verification (from \texttt{arm\_01\_inv\_grpo.ipynb}).}
\label{tab:invgrpo_rewards}
\begin{tabularx}{\textwidth}{l r r X}
\toprule
\textbf{Candidate Behavior} & \textbf{$\mathcal{R}$} & \textbf{$\hat{A}_i$} & \textbf{Policy Direction} \\
\midrule
Shortcut Mimic (copies $L_0$ onto $x'$)         & 1.00 & $-1.00$ & \textcolor{AlertRed}{Strongly suppressed} \\
Partial Solver (solves $x$ only, fails $x'$)    & 1.25 & $\approx 0.0$ & Neutral \\
Invariant Reasoner (solves both $x$ and $x'$)   & 2.50 & $+1.00$ & \textcolor{SuccessGreen}{Strongly reinforced} \\
\bottomrule
\end{tabularx}
\end{table}

\noindent Consistent with GRPO advantage formulation \cite{grpo_deepseek} and CLARity's consistency-aware RL \cite{clarity_rl}: invariant reasoning achieves $2.5\times$ the reward signal of shortcut mimicry, with clean normalized advantage separation.

%============================================================
% SECTION 4: FULL BUG POST-MORTEM
%============================================================
\section{Technical Bug Post-Mortem: All 9 Resolved Issues}
\label{sec:bugs}

\begin{table}[H]
\centering
\scriptsize
\caption{Complete Bug Resolution Record---Weeks 7 \& 8.}
\label{tab:bug_postmortem}
\begin{tabularx}{\textwidth}{l p{3.2cm} p{3.6cm} X}
\toprule
\textbf{ID} & \textbf{Symptom} & \textbf{Root Cause} & \textbf{Fix} \\
\midrule
BUG-1 & \texttt{SyntaxError} in sandbox   & Prompt injected as Python; \texttt{check()} called & Sanitized to \texttt{\# comments}; direct assertions. \\
BUG-2 & \texttt{TypeError: 'evaluation\_strategy'} & Removed in Transformers $\geq$ 4.41 & \texttt{eval\_strategy="no"} in \texttt{SFTConfig}. \\
BUG-3 & \texttt{TypeError: SFTTrainer 'tokenizer'} & TRL $\geq$ 0.24 moved to \texttt{SFTConfig} & \texttt{processing\_class=tokenizer}. \\
BUG-4 & \texttt{AttributeError: 'vanilla\_file'} & Wrong config namespace & Corrected to \texttt{settings.distillation}. \\
BUG-5 & \texttt{AttributeError: has no 'get'}  & Dataclass vs.\ dict mismatch & Added \texttt{.to\_dict()} + \texttt{\_normalize\_report()}. \\
BUG-6 & \texttt{FileNotFoundError: '../results'} & \texttt{os.chdir("..")} broke relative paths & All paths via \texttt{settings.storage}. \\
BUG-7 & \texttt{ReadTimeout: timeout=10}       & HF Hub 10\,s default on slow connection & \texttt{DEFAULT\_REQUEST\_TIMEOUT = 30}; local fallback. \\
BUG-8 & \texttt{IndentationError} in all eval  & Qwen generates body-only (no signature) \cite{qwen25coder} & \texttt{\_fix\_body\_indent()} in \texttt{sandbox.py}. \\
BUG-9 & GRPO \texttt{CUDA OOM} ($12.3$\,GB)   & Full $(G,T,V)$ tensor materialization \cite{grpo_deepseek} & Micro-batched \texttt{\_backward\_group\_loss()}; $3.03$\,GB. \\
\bottomrule
\end{tabularx}
\end{table}

%============================================================
% SECTION 5: VERIFIED ENVIRONMENT
%============================================================
\section{Verified Runtime Environment}

\begin{table}[H]
\centering
\small
\caption{Verified Hardware \& Software Stack (frozen in \texttt{requirements.txt} and \texttt{environment.yml}).}
\label{tab:env_specs}
\begin{tabularx}{\textwidth}{l l X}
\toprule
\textbf{Component} & \textbf{Version} & \textbf{Role} \\
\midrule
Python             & 3.10.20 (AMD64)           & Runtime interpreter \\
CUDA Toolkit       & 12.4                      & GPU kernel acceleration \\
Hardware           & NVIDIA RTX 3070 Ti (8\,GB GDDR6) & $<8.0$\,GB VRAM constraint \\
PyTorch            & 2.6.0+cu124               & Tensor \& autograd \\
Transformers       & 4.51.3                    & Causal LM \& tokenization \\
TRL                & 0.24.0                    & SFTTrainer, DPOTrainer, GRPOTrainer \\
PEFT               & 0.19.1                    & QLoRA adapters \cite{peft_qlora} \\
BitsAndBytes       & 0.49.2                    & 4-bit NF4 double quantization \\
Accelerate         & 1.7.0                     & Device mapping \& offloading \\
Datasets           & 3.6.0                     & Arrow streaming \& caching \\
\bottomrule
\end{tabularx}
\end{table}

%============================================================
% SECTION 6: FULL NOTEBOOK VERIFICATION TABLE
%============================================================
\section{Complete Notebook Verification Summary}

\begin{table}[H]
\centering
\small
\caption{End-to-End Execution Verification Across All 5 Core Notebooks + 4 Arm Notebooks.}
\label{tab:nb_verification}
\begin{tabularx}{\textwidth}{l l c X}
\toprule
\textbf{Notebook} & \textbf{Stage} & \textbf{Status} & \textbf{Key Verified Outputs} \\
\midrule
\texttt{nb\_01\_data\_pipeline}      & Stage 1: Acquisition  & \textcolor{SuccessGreen}{\textbf{PASS}} & 664 tasks cached ($L_0$: 164, $L_{1-5}$: 500) \\
\texttt{nb\_02\_baseline\_eval}      & Stage 2: Baseline     & \textcolor{SuccessGreen}{\textbf{PASS}} & M1: AUC 82.22\%, report 21.3\,MB \\
\texttt{nb\_03\_distillation}        & Stage 3: Distillation & \textcolor{SuccessGreen}{\textbf{PASS}} & 500 CoT samples, mean 377\,tok \\
\texttt{nb\_04\_qlora\_training}     & Stage 4: Training     & \textcolor{SuccessGreen}{\textbf{PASS}} & M2 checkpoint; peak VRAM 6.42\,GB \\
\texttt{nb\_05\_post\_training\_eval}& Stage 5: Eval         & \textcolor{SuccessGreen}{\textbf{PASS}} & M1 vs M2: 664 tasks; comparison plots \\
\addlinespace[3pt]
\texttt{arm\_01\_inv\_grpo}          & Arm 1: Inv-GRPO       & \textcolor{SuccessGreen}{\textbf{PASS}} & $\mathcal{R}_{inv}=2.50$; $\hat{A}=+1.00$; 3.03\,GB VRAM \\
\texttt{arm\_02\_contrastive\_sft}   & Arm 2: DPO            & \textcolor{SuccessGreen}{\textbf{PASS}} & DPO preference pairs; $y^+/y^-$ format \\
\texttt{arm\_03\_ast\_rl}            & Arm 3: AST-RL         & \textcolor{SuccessGreen}{\textbf{PASS}} & AST Jaccard: rename $>95\%$, wrong $<20\%$ \\
\texttt{arm\_04\_step\_rlvr}         & Arm 4: Step-RLVR      & \textcolor{SuccessGreen}{\textbf{PASS}} & Partial reward 0.70 vs.\ binary 0.00 \\
\bottomrule
\end{tabularx}
\end{table}

%============================================================
% SECTION 7: NEXT STEPS & OPEN QUESTION
%============================================================
\section{Week 9 Roadmap \& Open Research Questions}

\begin{keybox}[title={\faClock\ Week 9 Roadmap}]
\textbf{Priority 1 --- Full Inv-GRPO Training (Arm 1):}\\
Execute \texttt{arm\_01\_inv\_grpo.ipynb} to produce Model \textbf{M6} (Inv-GRPO). Target: recover M1 $L_0$ precision (93.29\%) while maintaining M2-level $L_{2+}$ robustness ($>90\%$). Theoretical prediction: \textbf{Ladder AUC $\geq$ 88\%} with \textbf{Memorization Risk $= 0.000$}.

\vspace{0.4em}
\textbf{Priority 2 --- Three-Model Comparative Evaluation:}\\
Run M6 through \texttt{nb\_05} and generate the definitive M1 vs.\ M2 vs.\ M6 comparison table and degradation curves.

\vspace{0.4em}
\textbf{Priority 3 --- Arm 2--4 Training Sweeps:}\\
Execute Contrastive DPO (Arm 2), AST-RL (Arm 3), and Step-RLVR (Arm 4) for the full 5-model ablation.
\end{keybox}

\begin{domainbox}[title={\faQuestion\ Open Scientific Question for Dr.\ Ghada Soliman}]
The M1 vs.\ M2 results reveal a stark trade-off: M2 achieves \textbf{zero memorization risk} but loses \textbf{28.7\,pp on $L_0$} canonical coding. The central hypothesis of this project is that Inv-GRPO resolves this by construction---the paired reward signal preserves both \emph{canonical fluency} (through $\mathcal{R}_{\text{exec}}(y_i)$ on the original prompt) and \emph{structural invariance} (through $\mathcal{R}_{\text{exec}}(y'_i)$ on the perturbed prompt).

\textbf{Key question:} Is Inv-GRPO a \textbf{strict Pareto improvement} over Vanilla SFT, or does a hybrid initialization (SFT M2 $\to$ Inv-GRPO) outperform pure RL from M1?
\end{domainbox}

%============================================================
% REFERENCES
%============================================================
\newpage
\printbibliography[heading=bibintoc, title={References \& Cited Literature}]

\end{document}
